\documentclass[11pt,reqno]{amsart}

\usepackage{amsmath,amssymb,amsthm}
\usepackage[T1]{fontenc}
\usepackage[expansion=false]{microtype}
\usepackage[margin=1in]{geometry}
\usepackage{booktabs}
\usepackage{hyperref}
\hypersetup{colorlinks=true,linkcolor=blue,citecolor=blue,urlcolor=blue}

\theoremstyle{plain}
\newtheorem{theorem}{Theorem}[section]
\newtheorem{proposition}[theorem]{Proposition}
\newtheorem{lemma}[theorem]{Lemma}
\newtheorem{corollary}[theorem]{Corollary}

\theoremstyle{definition}
\newtheorem{observation}[theorem]{Observation}

\theoremstyle{remark}
\newtheorem{remark}[theorem]{Remark}

\newcommand{\Q}{\mathbb{Q}}
\newcommand{\Z}{\mathbb{Z}}
\newcommand{\R}{\mathbb{R}}
\newcommand{\PP}{\mathbb{P}}
\newcommand{\HH}{\mathbb{H}}
\newcommand{\Aut}{\operatorname{Aut}}
\newcommand{\Bir}{\operatorname{Bir}}
\newcommand{\NS}{\operatorname{NS}}
\DeclareMathOperator{\rank}{rank}
\newcommand{\EPCP}{E_{\mathrm{PCP}}}

\begin{document}

\title[Structural and arithmetic obstructions for the perfect cuboid problem]
{Structural and arithmetic obstructions to elementary approaches to the perfect cuboid problem}

\author{C$\Lambda$ / Lightman Chang}
\address{Independent Researcher}
\email{lightman.chang@gmail.com}

\subjclass[2020]{Primary 11D09; Secondary 14J29, 14E07, 11G05, 11G50}

\keywords{perfect cuboid, Euler brick, surface of general type, birational automorphism, canonical height, multiplicative reduction}

\begin{abstract}
A \emph{perfect cuboid} is a rectangular box with integer edges, face diagonals,
and space diagonal; its existence is a long-standing open problem. Rational points
of the associated smooth complete intersection of four quadrics $V\subset\PP^6$
encode candidate cuboids. We record several structural facts about $V$ and about
the elliptic fibration of its Euler-brick K3 double cover, and use them to delimit
which elementary mechanisms can and cannot, in principle, contribute to a finiteness
statement. Our main observation is that no birational self-map of $V$ can effect a
height-strict infinite descent of the Markov--Vieta type: $V$ is a minimal surface
of general type (its canonical class equals the hyperplane class), so its birational
automorphism group is finite; a fixed Weil height is then bounded along any finite
birational orbit, and no strict per-step descent can persist. We compute this group
explicitly
($\Aut(V)=S_3\ltimes(\Z/2)^6$ of order $384$, all linear), record the surface
invariants $K_V^2=16$, $c_2=80$, $p_g=7$, $q=0$, and the Picard bounds
$16\le\rho_{\mathrm{geom}}(V')\le 20$ for the K3 cover. On the arithmetic side we
note that every fiber of the standard elliptic family has purely multiplicative
reduction, so all non-archimedean local heights are nonpositive, which obstructs an
elementary discriminant-free lower height bound; and we verify by descent and
Kolyvagin's theorem that the two rank-zero Jacobian factors of the genus-five fiber
are the curves $120a2$ and $80a1$. These are clarifying negatives that complement
the recent conditional and elementary work of Peschmann and of Yelle.
\end{abstract}

\maketitle

\section{Introduction}\label{sec:intro}

A \emph{perfect cuboid} (or perfect Euler brick) is a rectangular parallelepiped
with positive integer edges $a,b,c$, face diagonals $d,e,f$, and space diagonal
$g$, that is, a simultaneous integer solution of
\begin{equation}\label{eq:pcp}
a^2+b^2=d^2,\quad b^2+c^2=e^2,\quad a^2+c^2=f^2,\quad a^2+b^2+c^2=g^2 .
\end{equation}
Whether one exists is a classical open question; an extensive history is recorded
by Guy \cite{Guy}. The four conditions \eqref{eq:pcp} define a smooth complete
intersection of four quadrics
\[
V:\quad a^2+b^2-d^2=b^2+c^2-e^2=a^2+c^2-f^2=a^2+b^2+c^2-g^2=0
\]
in $\PP^6$, and a perfect cuboid is a $\Q$-point of $V$ with all coordinates
nonzero. Forgetting the fourth equation gives the Euler-brick surface
$V'\subset\PP^5$ cut out by the first three quadrics, classically known to be a K3
surface; $V\to V'$ is a double cover branched along $\{g=0\}$ \cite{vanLuijk}.

The problem has been studied recently from two complementary directions. Naskr\k{e}cki
\cite{Naskrecki} computed the geometric Mordell--Weil rank of the elliptic family
attached to the face conditions. Yoshida \cite{Yoshida} constructed, in the opposite
direction, families of \emph{face} cuboids. Peschmann \cite{Peschmann1,Peschmann2}
reduced the problem to a one-parameter family of genus-three hyperelliptic curves and
proved nonexistence on $1072$ explicit fibers by a torsion-intersection argument,
leaving genus-five Chabauty and uniform rank control as named open items. Yelle
\cite{Yelle} gave an elementary congruence/descent obstruction operating on
divisibility data along the space diagonal.

This note collects structural and arithmetic facts about $V$, its K3 cover $V'$, and
the standard per-fiber elliptic curve, with one purpose: to make precise which
\emph{elementary} mechanisms cannot, by their geometry alone, contribute to a
finiteness statement, and which auxiliary computations are settled unconditionally.
None of the results below resolves the perfect cuboid problem, and we make no such
claim. Their value is delimitative and complementary: they map the boundary between
methods that transfer to $V$ and methods that, for structural reasons, do not.

We single out one theorem.

\begin{theorem}\label{thm:main}
Let $V\subset\PP^6$ be the perfect cuboid complete intersection, a smooth minimal
surface of general type with ample canonical class $K_V=\mathcal{O}_V(1)$. Let $L$ be
any ample line bundle on $V$ and $H_L$ an associated Weil height. There is no
birational self-map $\sigma\colon V\dashrightarrow V$ and constant $C>0$ for which
\[
H_L(\sigma(P))\le H_L(P)-C
\]
holds for all $P$ in a $\sigma$-invariant Zariski-dense subset $S\subseteq V(\Q)$. In
particular, no Markov--Vieta or Fermat-style height-strict
infinite descent can be realized by a birational transformation of $V$.
\end{theorem}

Theorem~\ref{thm:main} explains, intrinsically, why the descent that succeeds for
single equations such as $x^4+y^4=z^2$ has no analogue acting on the cuboid surface
itself: the relevant ambient is a rigid surface of general type, not a curve, and its
birational symmetries form a finite group. It does not contradict the elementary
descent of Yelle \cite{Yelle}, whose mechanism is a prime-propagation argument on
divisibility data, not a birational self-map of $V$; the two operate on different
objects, and Section~\ref{sec:descent} makes the distinction explicit.

The remaining sections record: the explicit automorphism group and surface invariants
of $V$ and the Picard data of $V'$ (Section~\ref{sec:geometry}); the nonpositivity of
all non-archimedean local heights on the standard family and its consequence for
elementary height lower bounds (Section~\ref{sec:heights}); and the rank-zero
verification of the two distinguished Jacobian factors (Section~\ref{sec:ranks}).
Section~\ref{sec:remarks} comments honestly on multiplicative reformulations and on
the unlikeliness-of-intersection viewpoint, recording why neither supplies the
missing ingredient. All computations were performed in PARI/GP 2.15.4
\cite{PARI} and in SymPy; scripts and captured output accompany the paper
(Section~\ref{sec:repro}).

\section{The descent obstruction}\label{sec:descent}

We first prove Theorem~\ref{thm:main}; the geometric inputs are verified in
Section~\ref{sec:geometry}.

\begin{proof}[Proof of Theorem~\ref{thm:main}]
By adjunction for a smooth complete intersection of four quadrics in $\PP^6$, the
canonical class is
\[
K_V=\bigl(K_{\PP^6}+4\cdot\mathcal{O}(2)\bigr)\big|_V=(-7+8)\,\mathcal{O}_V(1)
   =\mathcal{O}_V(1),
\]
which is ample (Proposition~\ref{prop:invariants}). Hence $V$ is a minimal surface
of general type. By Matsumura's theorem \cite{Matsumura}, a smooth projective variety
with ample canonical bundle admits no positive-dimensional connected group of
birational transformations; together with the finiteness of the automorphism group of
a variety of general type (Maehara's extension of Severi's theorem
\cite{Maehara,Iitaka}), the group $\Bir(V)$ is finite. Moreover, for a minimal surface
of general type the natural map $\Aut(V)\to\Bir(V)$ is an isomorphism, since a
birational map between minimal models of a non-ruled surface is an isomorphism.

Let $\sigma\in\Bir(V)=\Aut(V)$. Since $\Bir(V)$ is finite, $\sigma$ has finite order
$k\ge1$, so $\sigma^{k}=\mathrm{id}$. Fix an ample $L$ and a representative Weil height
$H_L\colon V(\Q)\to\R$. Suppose, for contradiction, that $\sigma$ effected a
height-strict descent, that is,
\[
H_L(\sigma P)\le H_L(P)-C\qquad\text{for some constant }C>0\text{ and all }P\in S,
\]
where $S\subseteq V(\Q)$ is a $\sigma$-invariant Zariski-dense set. Summing this
inequality along the finite orbit $P,\sigma P,\dots,\sigma^{k-1}P$ of any $P\in S$ (each
$\sigma^{i}P\in S$ by $\sigma$-invariance, so the inequality applies at every step) gives
\[
H_L(\sigma^{k}P)\le H_L(P)-kC .
\]
But $\sigma^{k}=\mathrm{id}$, so the left-hand side equals $H_L(P)$, forcing $kC\le0$---a
contradiction, since $k\ge1$ and $C>0$. (Heights are well defined only up to $O(1)$; the
argument fixes one representative $H_L$ and evaluates it at the actual orbit points, where
the periodicity $\sigma^{k}P=P$ is exact, so the $O(1)$ ambiguity never enters.) Hence no
birational self-map of $V$ can effect a height-strict infinite descent on $V(\Q)$.
\end{proof}

\begin{remark}\label{rem:yelle}
The hypothesis is that the descent be effected \emph{by a birational self-map of}
$V$. This is exactly the structure used by the Markov tree for
$x^2+y^2+z^2=3xyz$ (the Vieta involution $z\mapsto 3xy-z$) and by the Fermat descent
for $x^4+y^4=z^2$ (where the ambient is a curve of genus three, not a general-type
surface). For the cuboid system, the Pythagorean relation $x^2+y^2=z^2$ is a
quadratic in $z$ with no linear cross term, so a single Berggren/Vieta step is only a
sign reflection; and a simultaneous step on the triples sharing a common edge imposes
an incidental linear constraint (for instance $b+f=c+d$ for the triples through $a$)
that a generic Euler brick fails---the classical brick $(44,117,240)$ has
$b+f=361\ne 365=c+d$. Theorem~\ref{thm:main} is the structural reason behind these
case-by-case failures. It says nothing about descents that are \emph{not}
birational self-maps of $V$; the divisibility descent of Yelle \cite{Yelle} is of that
other kind and is unaffected.
\end{remark}

\section{The automorphism group and surface invariants}\label{sec:geometry}

\begin{proposition}\label{prop:invariants}
The surface $V$ is a smooth minimal surface of general type with
\[
K_V=\mathcal{O}_V(1)\ \text{ample},\qquad K_V^2=16,\quad c_2(V)=80,\quad
\chi(\mathcal{O}_V)=8,\quad p_g=7,\quad q=0 .
\]
In particular $V$ is \emph{not} a K3 surface; only the Euler-brick double cover
$V'$ is K3.
\end{proposition}

\begin{proof}
With $H$ the hyperplane class, $c(T_{\PP^6})=(1+H)^7$ and the normal bundle of the
four quadrics contributes $c(N)=(1+2H)^4$, so modulo $H^3$
\[
c(T_V)\equiv(1+H)^7(1+2H)^{-4}\equiv 1-H+5H^2 .
\]
The degree of $V$ is $H^2\cdot V=2^4=16$. Reading off $c_1(T_V)=H$ gives
$K_V=H=\mathcal{O}_V(1)$, hence $K_V^2=16$ and $K_V$ ample. The second Chern number
is $c_2(V)=5\cdot 16=80$, so $\chi(\mathcal{O}_V)=(K_V^2+c_2)/12=8$. As a complete
intersection surface in $\PP^6$, the Lefschetz hyperplane theorem gives $q=0$, whence
$p_g=\chi-1=7$. Since $K_V$ is ample the Kodaira dimension is $2$ and $V$ is of
general type and minimal; a K3 surface has $K=0$, so $V$ is not K3. The symbolic
computation is in \texttt{03\_invariants\_general\_type.out}.
\end{proof}

\begin{proposition}\label{prop:aut}
The automorphism group of $V$ is
\[
\Aut(V)=\Bir(V)=S_3\ltimes(\Z/2)^6,\qquad |\Aut(V)|=384,
\]
and every element is the restriction of a linear automorphism of $\PP^6$. The $S_3$
factor permutes the edge labels $(a,b,c)$ with the induced permutation of the face
diagonals; the $(\Z/2)^6$ factor consists of the coordinate sign changes modulo the
global scalar.
\end{proposition}

\begin{proof}
Since $K_V=\mathcal{O}_V(1)$, the surface is canonically embedded, so every
automorphism is induced by a linear map of the ambient $\PP^6$ (an automorphism
preserving the canonical, hence the hyperplane, class is projectively linear).
A direct determination of the linear stabilizer of the defining ideal shows that its
Lie algebra is one-dimensional (only the scalar matrix), so the linear automorphism group
is $0$-dimensional, hence finite; an explicit enumeration of the stabilizer then gives
$S_3\ltimes(\Z/2)^6$ of order $384$: the six permutations of $(a,b,c)$ each lift to a
unique consistent permutation of $(d,e,f)$ fixing $g$, and the seven coordinate sign
flips, modulo the diagonal scalar in $\mathrm{PGL}_7$, give $(\Z/2)^6$ of order $64$. The
equality $\Aut(V)=\Bir(V)$ was established in the proof of Theorem~\ref{thm:main}. This
linear-stabilizer determination is a Gr\"obner-basis ideal computation carried out in the
project's prior \texttt{aut\_birV} package (\texttt{01\_linear\_aut.out},
\texttt{04\_s3\_verification.out}); it is not reproducible in PARI/GP, and we cite it as
an external computer-algebra input.
\end{proof}

The space diagonal satisfies $g^2=a^2+b^2+c^2$, a symmetric function of $(a,b,c)$, so
$g^2$ is fixed by all of $\Aut(V)$. The cuboid condition---rationality of $g$---is
therefore invariant under the full automorphism group, and no automorphism maps a
candidate cuboid to a degenerate or otherwise distinguished locus. This is the
geometric content behind Remark~\ref{rem:yelle}: the symmetry group reshuffles a
candidate among other candidates with the same space diagonal, never to a
contradiction.

The double cover is genuinely different.

\begin{proposition}\label{prop:k3}
The Euler-brick cover $V'\subset\PP^5$ (the first three quadrics) has twelve
$\Q$-rational ordinary double points; its minimal resolution $V'_{\min}$ is a K3
surface with geometric Picard rank $\rho_{\mathrm{geom}}(V'_{\min})\in[16,20]$.
\end{proposition}

\begin{proof}[Proof sketch]
The lower bound $\rho\ge 16$ comes from the explicit classes $H$, three fibration
classes $F_1,F_2,F_3$, and the twelve exceptional $(-2)$-curves over the nodes; their
$16\times16$ intersection matrix has full rank. The upper bound $\rho\le 20$ follows
from Frobenius-trace (van Luijk) computations at good primes $p=3,11,17$, where the
Artin parity and the counted Newton sums force the geometric Picard rank to be at
most $20$. See \cite{vanLuijk} for the method and the cited internal computation for
the trace data.
\end{proof}

Proposition~\ref{prop:k3} is the reason the rigidity of Theorem~\ref{thm:main} does
not extend to the cover: a K3 surface carrying $(-2)$-curves can have an infinite
birational automorphism group (an infinite reflection group), so the finite-group
argument is special to the general-type surface $V$ and cannot be invoked on $V'$.

\section{Multiplicative reduction and the height obstruction}\label{sec:heights}

The standard per-fiber elliptic curve attached to a Pythagorean parameter
$q=(m^2-n^2)/(2mn)$ is
\[
\EPCP(q):\quad y^2=x(x+1)(x+q^2)=x^3+(1+q^2)x^2+q^2x ,
\]
with full rational $2$-torsion. We record an arithmetic feature that bears on
elementary attempts to bound canonical heights from below.

\begin{proposition}\label{prop:mult}
For every Pythagorean parameter $q$, the curve $\EPCP(q)$ has purely multiplicative
(Kodaira type $I_n$) reduction at every prime of bad reduction. Consequently every
non-archimedean N\'eron local height is nonpositive:
\[
\lambda_p(P)=-\frac{n_c(N_p-n_c)}{N_p}\log p\ \in\ \Bigl[-\tfrac{N_p}{4}\log p,\,0\Bigr],
\qquad N_p=v_p(\Delta_{\min}),
\]
where $n_c\in[0,N_p/2]$ is the component index. Summing over bad primes,
$\sum_p\lambda_p(P)\in[-\tfrac14\log|\Delta_{\min}|,\,0]$.
\end{proposition}

\begin{proof}
Multiplicativity is the statement $v_p(c_4)=0$ at every bad prime, which we verify on
representative fibers in \texttt{02\_local\_heights\_In.out}; for example at
$(m,n)=(11,2)$ the bad primes are $3,7,11,13,23,73$, each with $v_p(c_4)=0$ and
Kodaira type $I_n$. The local-height formula for multiplicative reduction is N\'eron's
\cite[Ch.~VI]{SilvermanATAEC}; its value is a nonpositive multiple of $\log p$, with
extremes $0$ on the identity component and $-\tfrac{N_p}{4}\log p$ at the deepest
component. The sum bound follows from $\sum_p N_p\log p=\log|\Delta_{\min}|$.
\end{proof}

\begin{observation}\label{obs:noabs}
On this family the positive part of the canonical height
$\widehat{h}(P)=\lambda_\infty(P)+\sum_p\lambda_p(P)$ resides entirely in the
archimedean term $\lambda_\infty$, since the non-archimedean total is nonpositive by
Proposition~\ref{prop:mult}. The archimedean local height is governed by the elliptic
logarithm and has no lower bound that grows with $\log|\Delta_{\min}|$; numerically the
ratio $\widehat{h}_{\min}/\log|\Delta_{\min}|$ over a representative sample of
rank-positive fibers decreases as the Szpiro ratio grows (the detailed sweep is part of
the project's companion Szpiro-ratio analysis). Thus no elementary,
discriminant-only lower bound $\widehat{h}(P)\ge c\log|\Delta_{\min}|$ with $c$
independent of the Szpiro ratio is visible for $\EPCP(q)$: the multiplicative
structure supplies no positive height budget. A Szpiro-uniform lower bound of
Lang--Silverman type is available in principle but is non-effective \cite{Wagener},
and the effective per-fiber bound of Petsche \cite{Petsche} degrades polynomially in
the Szpiro ratio.
\end{observation}

Observation~\ref{obs:noabs} is delimitative: it identifies the precise reason an
elementary height inequality cannot be the missing input, and it locates the only
remaining elementary handle in a bound on the Szpiro ratio rather than in the height
directly.

\section{The rank-zero Jacobian factors}\label{sec:ranks}

The genus-five fiber arising from the face fibration of $V$ has a Jacobian that
decomposes, up to $\Q$-isogeny, into three elliptic factors of rank one and two
further elliptic factors $X_+,X_-$ \cite[\S2]{vanLuijk}. We record the unconditional
status of the latter two.

\begin{proposition}\label{prop:xpm}
The two distinguished elliptic factors are, in Cremona labels,
\[
X_+=120a2:\ y^2=x^3+x^2-20x,\qquad X_-=80a1:\ y^2=x^3-7x+6,
\]
each of analytic rank $0$ with $L(X_\pm,1)>0$. By Kolyvagin's theorem
\cite{Kolyvagin} both have Mordell--Weil rank $0$ over $\Q$, unconditionally; thus
$X_+(\Q)$ and $X_-(\Q)$ are finite.
\end{proposition}

\begin{proof}
We compute the analytic ranks, the special values
$L(X_+,1)=1.26949\ldots$ and $L(X_-,1)=1.00945\ldots$, the torsion subgroups
($\Z/4\times\Z/2$ for $X_+$, $(\Z/2)^2$ for $X_-$), and a $2$-descent rank bound of
$0$ in \texttt{01\_xpm\_kolyvagin.out}; the two integral models recorded in the
internal corpus have equal $j$-invariants and are confirmed to be these Cremona
curves by \texttt{ellidentify}. Since the analytic rank is $0$ and the central
$L$-value is nonzero, Kolyvagin's Euler-system argument \cite{Kolyvagin}, building on
Gross--Zagier \cite{GrossZagier}, gives algebraic rank $0$.
\end{proof}

Proposition~\ref{prop:xpm} is a clean ingredient---it removes two of the five
Jacobian factors from any rank count---rather than a result in its own right; the
finiteness it supplies is the routine consequence of a nonvanishing central $L$-value.

\section{Approaches that do not transfer}\label{sec:remarks}

We record two further negatives, stated to spare future effort rather than to
refute any published claim.

\begin{remark}[Multiplicative reformulations]\label{rem:mult}
Encoding the Pythagorean relations in $\Z[i]$, or the four-square relation in the
Hurwitz quaternions or the imaginary octonions, reproduces the cuboid system without
adding an obstruction. Each face relation $x^2+y^2=z^2$ factors in $\Z[i]$ as a
Gaussian square up to a squarefree inert part, and one recovers the classical fact
that a perfect cuboid forces each of $d,e,f,g$ to be composite (a face diagonal that
is prime is impossible)---a restatement of a Pocklington-type divisibility lemma. But
the shared edges $a,b,c$ couple the factorizations \emph{over} $\Z$, not over $\Z[i]$,
and any multiplicative identity derived from \eqref{eq:pcp} (for instance
$(a^2+bc)^2+a^2(b-c)^2=d^2f^2$) is an automatic consequence of those equations. The
non-commutative and non-associative settings add the parity normalization and norm
identities but no entanglement at the ideal level: the obstruction, if any, is a
rational-point question on the surface, not a unique-factorization question. We are
not aware of any claim in the literature that a purely multiplicative method resolves
the problem; this remark records why such a method cannot.
\end{remark}

\begin{remark}[Unlikely intersections]\label{rem:pz}
The $q$-line is the genus-zero modular curve $X(\Gamma_1(4)\cap\Gamma(2))$, with $q$
a Hauptmodul and a degree-twelve map to the $j$-line,
$j(q)=256(q^4-q^2+1)^3/(q^4(q^2-1)^2)$. One can frame the cuboid condition as an
intersection in the relative square of the universal elliptic curve over this base, a
mixed Shimura threefold whose modular shadow is $X_1(4)\times X_1(4)$. A dimension
count places a hypothetical one-parameter family of cuboids in excess dimension, hence
``atypical''; but an isolated perfect cuboid sits at the \emph{expected} dimension
zero and is not atypical, and the body-diagonal condition $1+q^2+c^2\in(\Q^\times)^2$
is a quadratic-twist squareness, not a special subvariety of the Shimura variety.
The Pila--Zannier method \cite{PilaZannier} and the atypical-intersection theorem of
Habegger--Pila \cite{HabeggerPila} therefore have no special locus to act on for the
fourth condition. We do not present this as a defeat of any proposed proof---no such
proof has been advanced in the literature---but as a precise statement of the
ingredient (a modular or height-growth structure for the body-diagonal condition)
that the unlikely-intersection viewpoint would require and that is presently absent.
\end{remark}

\section{Reproducibility}\label{sec:repro}

Computations use PARI/GP 2.15.4 \cite{PARI} and SymPy on a single x86-64 core. The
directory \texttt{scripts/} contains, with captured output:
\texttt{01\_xpm\_kolyvagin.gp} (Proposition~\ref{prop:xpm}: conductors, analytic
ranks, $L$-values, torsion, \texttt{ellrank}, \texttt{ellidentify} for both integral
models of $X_\pm$); \texttt{02\_local\_heights\_In.gp}
(Proposition~\ref{prop:mult}: multiplicative reduction and the local-height range at
each bad prime on three sample fibers); and
\texttt{03\_invariants\_general\_type.py} (Proposition~\ref{prop:invariants}: the
Chern computation $K_V^2=16$, $c_2=80$, $\chi=8$, $p_g=7$, $q=0$, and the group
order $384$). The automorphism-group symbolic verifications
(\texttt{01\_linear\_aut}, \texttt{04\_s3\_verification}) are part of the project's
prior \texttt{aut\_birV} computation referenced in Proposition~\ref{prop:aut}.

\begin{thebibliography}{99}

\bibitem{GrossZagier}
B.~H. Gross and D.~B. Zagier,
\emph{Heegner points and derivatives of $L$-series},
Invent. Math. \textbf{84} (1986), 225--320.

\bibitem{Guy}
R.~K. Guy,
\emph{Unsolved Problems in Number Theory}, 3rd ed.,
Springer, New York, 2004.

\bibitem{HabeggerPila}
P.~Habegger and J.~Pila,
\emph{O-minimality and certain atypical intersections},
Ann. Sci. \'Ec. Norm. Sup\'er. (4) \textbf{49} (2016), 813--858.

\bibitem{HindrySilverman}
M.~Hindry and J.~H. Silverman,
\emph{Diophantine Geometry: An Introduction},
Graduate Texts in Mathematics \textbf{201}, Springer, New York, 2000.

\bibitem{Iitaka}
S.~Iitaka,
\emph{Algebraic Geometry: An Introduction to Birational Geometry of Algebraic
Varieties},
Graduate Texts in Mathematics \textbf{76}, Springer, New York, 1982.

\bibitem{Kolyvagin}
V.~A. Kolyvagin,
\emph{Finiteness of $E(\mathbb{Q})$ and the Shafarevich--Tate group for a subclass
of Weil curves},
Izv. Akad. Nauk SSSR Ser. Mat. \textbf{52} (1988), 522--540;
English transl. Math. USSR-Izv. \textbf{32} (1989), 523--541.

\bibitem{Maehara}
K.~Maehara,
\emph{A finiteness property of varieties of general type},
Math. Ann. \textbf{262} (1983), 101--123.

\bibitem{Matsumura}
H.~Matsumura,
\emph{On algebraic groups of birational transformations},
Atti Accad. Naz. Lincei Rend. Cl. Sci. Fis. Mat. Natur. (8) \textbf{34} (1963),
151--155.

\bibitem{Naskrecki}
B.~Naskr\k{e}cki,
\emph{Mordell--Weil ranks of families of elliptic curves associated to Pythagorean
triples},
Acta Arith. \textbf{160} (2013), 159--183.

\bibitem{PARI}
The PARI~Group,
\emph{PARI/GP version 2.15.4},
Univ. Bordeaux, 2023, \url{https://pari.math.u-bordeaux.fr/}.

\bibitem{Peschmann1}
R.~Peschmann,
\emph{Quartic reductions and elliptic obstructions for perfect Euler bricks},
preprint, arXiv:2604.09328 (2026).

\bibitem{Peschmann2}
R.~Peschmann,
\emph{A torsion-intersection proof of perfect-cuboid nonexistence on $1072$ explicit
master-tuple fibers},
preprint, arXiv:2604.28072 (2026).

\bibitem{Petsche}
C.~Petsche,
\emph{Small rational points on elliptic curves over number fields},
New York J. Math. \textbf{12} (2006), 257--268.

\bibitem{PilaZannier}
J.~Pila and U.~Zannier,
\emph{Rational points in periodic analytic sets and the Manin--Mumford conjecture},
Atti Accad. Naz. Lincei Rend. Lincei Mat. Appl. \textbf{19} (2008), 149--162.

\bibitem{SilvermanATAEC}
J.~H. Silverman,
\emph{Advanced Topics in the Arithmetic of Elliptic Curves},
Graduate Texts in Mathematics \textbf{151}, Springer, New York, 1994.

\bibitem{vanLuijk}
R.~van Luijk,
\emph{On perfect cuboids},
Master's thesis, Universiteit Utrecht, 2000.

\bibitem{Wagener}
B.~Wagener,
\emph{About the proof of Lang's height conjecture},
preprint, arXiv:1706.03622 (2017).

\bibitem{Yelle}
S.~Yelle,
\emph{An elementary obstruction to the existence of a perfect cuboid},
preprint, arXiv:2602.00239 (2026).

\bibitem{Yoshida}
T.~Yoshida,
\emph{The relationship between face cuboids and elliptic curves},
preprint, arXiv:2407.09825 [math.NT] (2024).

\end{thebibliography}

\end{document}
